\documentclass[../../main.tex]{subfiles}

\graphicspath{{../../image/}{../../figure/}}

\begin{document}

\section{熵--熵流对}\label{section:ch03-熵流对}

设平衡律方程组\eqref{align:ch01-closed-system}%

（现写为\eqref{align:ch03-平衡律一般形式}）配有一个伴随平衡律\eqref{align:ch01-companion-law}%

。令 $Q_{k}\equiv\eta$，并把 $Q$ 重新指派为 $m$ 维行向量 $(Q_{1},\ldots,Q_{m})$，于是\eqref{align:ch01-companion-law}%

在新记号下写为
\begin{align}
\partial_{t}\eta(U(x,t),x,t)+\operatorname{div}Q(U(x,t),x,t)=h(U(x,t),x,t).
\label{align:ch03-熵平衡律}
\end{align}

正如我们将在\cref{section:ch03-例子}%

看到的，在连续介质物理的应用中，形如\eqref{align:ch03-熵平衡律}的伴随平衡律与热力学第二定律密切相关。因此，$\eta$ 称为平衡律方程组\eqref{align:ch03-平衡律一般形式}的\emph{熵}，而 $Q$ 称为与 $\eta$ 关联的\emph{熵流}。

\eqref{align:ch01-companion}%

当 $\alpha=k$ 时应写为
\begin{align}
D\eta(U,x,t)=B(U,x,t)^{\top}DH(U,x,t).
\label{align:ch03-Deta-DH}
\end{align}

设方程组已取正则形式\eqref{align:ch03-正则形式}，于是\eqref{align:ch03-Deta-DH} 简化为 $D\eta=B^{\top}$。此时\eqref{align:ch01-companion}%

与可积性条件\eqref{align:ch01-integrability}%

成为
\begin{align}
DQ_{\alpha}(U,x,t)=D\eta(U,x,t)\,DG_{\alpha}(U,x,t),\qquad \alpha=1,\cdots,m,
\label{align:ch03-DQ}
\end{align}
\begin{align}
D^{2}\eta(U,x,t)\,DG_{\alpha}(U,x,t)=DG_{\alpha}(U,x,t)^{\top}D^{2}\eta(U,x,t),\qquad \alpha=1,\cdots,m.
\label{align:ch03-可积性}
\end{align}

注意\eqref{align:ch03-可积性} 对单个未知函数 $\eta$ 施加了 $\frac{1}{2}n(n-1)m$ 个条件。因此，如\cref{section:ch01-伴随平衡律}%

已指出的那样，为\eqref{align:ch03-平衡律一般形式} 确定一个非平凡熵--熵流对的问题在形式上是超定的，除非 $n=1$ 而 $m$ 任意，或 $n=2$ 且 $m=1$。然而，当方程组对称时，可以取 $\eta=\frac{1}{2}|U|^{2}$ 使\eqref{align:ch03-可积性} 成立。反之，若\eqref{align:ch03-可积性} 成立且 $\eta(U,x,t)$ 关于 $U$ 一致凸，则状态向量的变换 $U^{*}=D\eta(U,x,t)^{\top}$ 使方程组成为对称的。于是，配有凸熵的正则形式\eqref{align:ch03-正则形式} 平衡律方程组必然是双曲的。

可积性条件还有一个有趣的等价形式：把\eqref{align:ch03-可积性} 沿任意方向 $\nu\in\mathbb{S}^{m-1}$ 投影，再分别左乘 $R_{j}(\nu;U,x,t)^{\top}$、右乘 $R_{k}(\nu;U,x,t)$（$j\neq k$）。只要 $\lambda_{j}(\nu;U,x,t)\neq\lambda_{k}(\nu;U,x,t)$，这一计算给出
\begin{align}
R_{j}(\nu;U,x,t)^{\top}D^{2}\eta(U,x,t)R_{k}(\nu;U,x,t)=0,\qquad j\neq k.
\label{align:ch03-可积性投影}
\end{align}

此外，即使 $\lambda_{j}(\nu;U,x,t)=\lambda_{k}(\nu;U,x,t)$，只要在该多重特征值的特征子空间中适当选取特征向量 $R_{j}(\nu;U,x,t)$ 与 $R_{k}(\nu;U,x,t)$，\eqref{align:ch03-可积性投影} 仍然成立。

注意\eqref{align:ch03-可积性投影} 对每一固定的 $\nu$ 都对 $\eta$ 施加了 $\frac{1}{2}n(n-1)$ 个条件，从而对 $\mathbb{S}^{m-1}$ 中 $m$ 个线性无关方向上的 $\nu$（由此覆盖全体 $\nu$）总共施加了 $\frac{1}{2}n(n-1)m$ 个条件。一个值得注意的例外是：方程组中各通量分量的 Jacobi 矩阵两两可交换，即
\begin{align}
DG_{\alpha}(U,x,t)\,DG_{\beta}(U,x,t)=DG_{\beta}(U,x,t)\,DG_{\alpha}(U,x,t),\qquad \alpha,\beta=1,\ldots,m.
\label{align:ch03-Jacobi可交换}
\end{align}
事实上，此时 $R_{i}(\nu;U,x,t)$ 不随 $\nu$ 变化，因而\eqref{align:ch03-可积性投影} 对 $\eta$ 只施加 $\frac{1}{2}n(n-1)$ 个条件。我们将在\cref{section:ch06-方程组L1理论}%

重新讨论这一非常特殊的方程组类。

一维空间维数下\eqref{align:ch03-可积性投影} 的超定性问题将在\cref{section:ch07-熵熵流对}%

深入考察。

\end{document}
